% modul-4.tex - Modul 4: Agentic Coding

\startcomponent modul-4
\environment ../pyt-layout
\environment ../pyt-env

\chapter{Agentic Coding -- KI-gestützte Entwicklung}

{\it Dauer: 1 Tag (4 Lektionen à 50 Minuten) • Voraussetzung: Module 1-3 abgeschlossen}

\startlernziele
Nach diesem Modul können Sie:
\startitemize[n,packed]
    \item KI als Entwicklungspartner effektiv nutzen
    \item Komplexe Prompts für Code-Generierung schreiben
    \item Test-Driven Development (TDD) mit KI praktizieren
    \item Code-Reviews mit KI durchführen
    \item Refactoring mit KI-Unterstützung
    \item Dokumentation automatisch generieren
    \item Grenzen und Best Practices verstehen
\stopitemize
\stoplernziele

% =============================================================================
% VORBEREITUNG
% =============================================================================

\section{Vorbereitung}

{\it Zeitaufwand: 2-3 Stunden vor dem Präsenzunterricht}

\subsection{Leseauftrag: Agentic Coding Konzepte}

\startinfobox
{\bf Was ist Agentic Coding?}

Entwicklung, bei der KI als aktiver Partner agiert, nicht nur als Werkzeug. Die KI generiert, erklärt, verbessert und dokumentiert Code -- der Mensch trifft die Entscheidungen.
\stopinfobox

\subsubsection{Die 4 Prinzipien}

\startitemize[n,packed]
    \item {\bf Iterative Verfeinerung} -- Prompts schrittweise verbessern
    \item {\bf Kontext ist König} -- Mehr Kontext = bessere Ergebnisse
    \item {\bf Verify \& Validate} -- Immer Code verstehen und testen
    \item {\bf Human in the Loop} -- Entwickler trifft finale Entscheidungen
\stopitemize

\subsection{KI-Tool-Vergleich}

Experimentieren Sie mit verschiedenen KI-Tools:

\starttabulate[|l|l|p(5cm)|]
\HL
\NC {\bf Tool} \NC {\bf Stärke} \NC {\bf Einsatz} \NC\NR
\HL
\NC ChatGPT \NC Erklärungen \NC Konzepte verstehen, Debugging \NC\NR
\NC Claude \NC Lange Kontexte \NC Grosse Codebasen analysieren \NC\NR
\NC GitHub Copilot \NC IDE-Integration \NC Code-Vervollständigung \NC\NR
\NC Cursor \NC Agentic Features \NC Komplexe Refactorings \NC\NR
\HL
\stoptabulate

% =============================================================================
% PRAXIS
% =============================================================================

\section{Praxis}

{\it Präsenzunterricht: 4 Lektionen à 50 Minuten}

\subsection{Lektion 1: Advanced Prompt Engineering}

{\it Dauer: 50 Minuten}

\subsubsection{Prompt Engineering Patterns}

{\bf 1. Role-Based Prompting:}
\startbashcode
Du bist ein erfahrener Python-Entwickler mit Fokus auf
Clean Code und Test-Driven Development.
\stopbashcode

{\bf 2. Context-Rich Prompting:}
\startbashcode
Gegeben: [Kontext - bestehender Code, Anforderungen]
Aufgabe: [Was soll gemacht werden]
Constraints: [Einschraenkungen, Stil-Vorgaben]
\stopbashcode

{\bf 3. Chain-of-Thought:}
\startbashcode
Erklaere Schritt fuer Schritt, wie du die Loesung entwickelst.
Zeige deine Ueberlegungen.
\stopbashcode

{\bf 4. Few-Shot Learning:}
\startbashcode
Hier sind 2 Beispiele fuer die gewuenschte Funktion:
[Beispiel 1]
[Beispiel 2]
Jetzt mache das Gleiche fuer: [neue Anforderung]
\stopbashcode

\subsubsection{Prompt-Template}

\startbashcode
[Rolle]: Du bist ein erfahrener Python-Entwickler.

[Kontext]: Ich arbeite an einer Web-API mit FastAPI.
Der bestehende Code verwendet Pydantic-Models.

[Aufgabe]: Erstelle einen Endpoint fuer User-Registrierung.

[Anforderungen]:
- Email-Validierung
- Passwort-Hashing mit bcrypt
- Duplikat-Pruefung
- Aussagekraeftige Fehlermeldungen

[Beispiel]:
Eingabe: {"email": "test@example.com", "password": "Secure123!"}
Ausgabe: {"id": 1, "email": "test@example.com", "created_at": "..."}

[Constraints]:
- Type Hints verwenden
- Docstrings auf Deutsch
- Keine print-Statements
\stopbashcode

\subsection{Lektion 2: Test-Driven Development mit KI}

{\it Dauer: 50 Minuten}

\subsubsection{TDD-Workflow}

\startitemize[n,packed]
    \item {\bf Red:} Test schreiben, der fehlschlägt
    \item {\bf Green:} Minimalen Code schreiben, der Test besteht
    \item {\bf Refactor:} Code verbessern, Tests müssen weiter bestehen
\stopitemize

\subsubsection{Tests zuerst mit KI}

\startpythoncode
# Prompt: "Schreibe pytest-Tests fuer eine Funktion
# die Primzahlen prueft"

import pytest
from primzahl import ist_primzahl

def test_zwei_ist_primzahl():
    assert ist_primzahl(2) == True

def test_eins_ist_keine_primzahl():
    assert ist_primzahl(1) == False

def test_negative_zahl():
    assert ist_primzahl(-5) == False

def test_grosse_primzahl():
    assert ist_primzahl(97) == True

def test_grosse_nicht_primzahl():
    assert ist_primzahl(100) == False
\stoppythoncode

\subsubsection{Implementation generieren}

\startpythoncode
# Prompt: "Implementiere ist_primzahl() so dass
# alle Tests bestehen"

def ist_primzahl(n: int) -> bool:
    """Prueft, ob n eine Primzahl ist."""
    if n < 2:
        return False
    if n == 2:
        return True
    if n % 2 == 0:
        return False
    for i in range(3, int(n ** 0.5) + 1, 2):
        if n % i == 0:
            return False
    return True
\stoppythoncode


\subsection{Lektion 3: Code Review \& Refactoring}

{\it Dauer: 50 Minuten}

\subsubsection{KI als Code Reviewer}

\startbashcode
Analysiere folgenden Code auf:
- Code Smells
- Moegliche Bugs
- Performance-Probleme
- Sicherheitsluecken
- Verbesserungsvorschlaege

[Code hier einfuegen]
\stopbashcode

\subsubsection{Refactoring-Beispiel}

{\bf Vorher:}
\startpythoncode
def f(x, y):
    r = []
    for i in x:
        if i > y:
            r.append(i * 2)
    return r
\stoppythoncode

{\bf Nachher (nach KI-Refactoring):}
\startpythoncode
def filter_and_double(
    numbers: list[int],
    threshold: int
) -> list[int]:
    """
    Filtert Zahlen groesser als threshold
    und verdoppelt sie.
    """
    return [n * 2 for n in numbers if n > threshold]
\stoppythoncode

\subsubsection{Refactoring-Checkliste}

\startitemize[packed]
    \item Aussagekräftige Namen für Variablen/Funktionen
    \item Type Hints hinzufügen
    \item Docstrings schreiben
    \item Komplexe Logik vereinfachen
    \item Duplikate entfernen (DRY)
    \item Magic Numbers durch Konstanten ersetzen
\stopitemize

\subsection{Lektion 4: Dokumentation \& Best Practices}

{\it Dauer: 50 Minuten}

\subsubsection{Docstrings generieren}

\startbashcode
Generiere Google-Style Docstrings fuer alle
Funktionen in folgendem Code:
[Code]
\stopbashcode

\startpythoncode
def berechne_rabatt(
    preis: float,
    rabatt_prozent: float
) -> float:
    """Berechnet den rabattierten Preis.

    Args:
        preis: Der urspruengliche Preis in CHF.
        rabatt_prozent: Der Rabatt in Prozent (0-100).

    Returns:
        Der reduzierte Preis nach Abzug des Rabatts.

    Raises:
        ValueError: Wenn rabatt_prozent < 0 oder > 100.

    Example:
        >>> berechne_rabatt(100.0, 20.0)
        80.0
    """
    if not 0 <= rabatt_prozent <= 100:
        raise ValueError("Rabatt muss zwischen 0 und 100 sein")
    return preis * (1 - rabatt_prozent / 100)
\stoppythoncode

\subsubsection{README generieren}

\startbashcode
Erstelle eine README.md fuer folgendes Projekt:
- Projektname: [Name]
- Beschreibung: [Beschreibung]
- Hauptfeatures: [Features]
- Struktur: [Dateien/Ordner]

Verwende Badges, Inhaltsverzeichnis,
Installation, Verwendung, Beispiele.
\stopbashcode

% =============================================================================
% ÜBUNGEN
% =============================================================================

\section{Übungen}

\startaufgabenbox
{\bf Übung 1: Komplexe Feature-Implementierung} {\it (20 Min.)}

Mit KI-Unterstützung implementieren:
\startitemize[packed]
    \item URL-Shortener mit Statistiken
    \item Passwort-Generator mit Konfiguration
    \item Markdown-zu-HTML-Konverter
\stopitemize

Dokumentieren Sie Ihren Prompt-Iterationsprozess!
\stopaufgabenbox

\blank[medium]

\startaufgabenbox
{\bf Übung 2: TDD-Challenge} {\it (20 Min.)}

TDD-Workflow durchführen:
\startitemize[n,packed]
    \item Tests für Taschenrechner schreiben lassen
    \item Implementation generieren
    \item Edge Cases ergänzen
    \item Refactoring durchführen
\stopitemize
\stopaufgabenbox

\blank[medium]

\startaufgabenbox
{\bf Übung 3: Legacy Code Refactoring} {\it (15 Min.)}

Schlechten Code analysieren und verbessern:
\startitemize[packed]
    \item Code Smells identifizieren
    \item Verbesserungen vorschlagen lassen
    \item Refactoring durchführen
    \item Tests ergänzen
\stopitemize
\stopaufgabenbox

\blank[medium]

\startaufgabenbox
{\bf Übung 4: Vollständige Dokumentation} {\it (15 Min.)}

Für ein Mini-Projekt erstellen:
\startitemize[packed]
    \item Docstrings für alle Funktionen
    \item README.md mit Installation
    \item CHANGELOG.md
    \item Inline-Kommentare
\stopitemize
\stopaufgabenbox

% =============================================================================
% NACHBEARBEITUNG
% =============================================================================

\section{Nachbearbeitung}

{\it Hausaufgaben: 6-8 Stunden vor Modul 5}

\subsection{Aufgabe 1: Vollständiges CLI-Tool}

Entwickeln Sie ein CLI-Tool mit KI-Unterstützung:

\startbashcode
python todo-cli.py add "Einkaufen gehen"
python todo-cli.py list
python todo-cli.py done 1
python todo-cli.py delete 1
\stopbashcode

\startitemize[packed]
    \item TDD-Workflow anwenden
    \item Vollständige Dokumentation
    \item Fehlerbehandlung
    \item Persistenz (JSON-Datei)
\stopitemize

\subsection{Aufgabe 2: Code-Qualitäts-Analyse}

Analysieren Sie ein Open-Source-Projekt:
\startitemize[packed]
    \item 3 Code Smells identifizieren
    \item Verbesserungsvorschläge dokumentieren
    \item Refactoring-Plan erstellen
\stopitemize

\subsection{Aufgabe 3: Portfolio-Projekt}

Beginnen Sie Ihr Abschlussprojekt:
\startitemize[packed]
    \item Projektidee definieren
    \item README.md erstellen
    \item Grundstruktur mit KI entwickeln
    \item Erste Tests schreiben
\stopitemize

% =============================================================================
% MATERIALIEN
% =============================================================================

\section{Materialien}

\subsection{Code Review Checkliste}

\startitemize[packed]
    \item[$\square$] Funktioniert der Code wie erwartet?
    \item[$\square$] Sind alle Edge Cases behandelt?
    \item[$\square$] Gibt es ausreichend Tests?
    \item[$\square$] Sind Variablennamen aussagekräftig?
    \item[$\square$] Ist der Code DRY (keine Duplikate)?
    \item[$\square$] Sind Funktionen klein und fokussiert?
    \item[$\square$] Gibt es Docstrings/Kommentare?
    \item[$\square$] Werden Fehler sinnvoll behandelt?
    \item[$\square$] Gibt es Sicherheitsrisiken?
    \item[$\square$] Ist der Code performant?
\stopitemize

\subsection{Prompt Engineering Patterns}

\starttabulate[|l|p(8cm)|]
\HL
\NC {\bf Pattern} \NC {\bf Beschreibung} \NC\NR
\HL
\NC Role-Based \NC KI eine Expertenrolle zuweisen \NC\NR
\NC Context-Rich \NC Viel Kontext und Constraints geben \NC\NR
\NC Chain-of-Thought \NC Schrittweise Erklärung verlangen \NC\NR
\NC Few-Shot \NC Beispiele zeigen, dann generalisieren \NC\NR
\NC Iterative \NC Schrittweise verfeinern \NC\NR
\HL
\stoptabulate

\blank[big]

\starttippbox
{\bf Erfolgskriterien für Modul 4}

Sie haben das Modul erfolgreich abgeschlossen, wenn Sie:
\startitemize[packed]
    \item Komplexe Prompts für Code-Generierung schreiben können
    \item TDD-Workflow mit KI-Unterstützung durchführen können
    \item Code-Reviews mit KI durchführen können
    \item Vollständige Dokumentation erstellen können
    \item Die Grenzen von KI-generiertem Code verstehen
\stopitemize
\stoptippbox

\stopcomponent

